\section{Metodologi Penelitian}

Penelitian ini menggunakan pendekatan \textit{Design Science Research Methodology} (DSRM)\label{first:dsrm}, yaitu paradigma yang berorientasi pada perancangan dan evaluasi artefak untuk menyelesaikan permasalahan yang teridentifikasi, alih-alih sekadar menjelaskan atau memprediksi fenomena yang telah berlangsung \parencite{peffers_design_2007}.
Titik masuk yang diikuti adalah \textit{problem-centered initiation}, karena penelitian ini bermula dari kesenjangan yang teridentifikasi pada tinjauan pustaka.
DSRM menjadi payung paradigma, sedangkan aktivitas \textit{Demonstration} dan \textit{Evaluation} dioperasionalkan melalui desain eksperimen kuantitatif terkontrol.
Jadi artefak akan dibandingkan dengan enam algoritma \textit{baseline} pada dataset yang dikonstruksi, diukur dengan metrik yang ditetapkan, dan diuji signifikansinya secara statistik.
Proses penelitian mengikuti enam aktivitas DSRM \parencite{peffers_design_2007} berikut:

\begin{enumerate}
    \item \textit{Problem identification and motivation} -- Mengidentifikasi kesenjangan penelitian, yaitu belum tersedianya algoritma metaheuristik untuk \textit{policy mining} ABAC yang sekaligus ringan, sadar dependensi, dan dapat diperbarui inkremental pada perangkat \textit{edge-IoT} beserta urgensinya.
    \item \textit{Define objectives of a solution} -- Menetapkan empat tujuan khusus beserta kriteria keberhasilan terukur yang dirumuskan sebagai Hipotesis H1-H4.
    \item \textit{Design and development} -- Merancang dan mengimplementasikan artefak penelitian. Aktivitas ini paling ekstensif dan mencakup tiga komponen: (a) perancangan algoritma CoDA-EDA itu sendiri -- representasi model probabilistik, mekanisme pembaruan kompetitif tanpa populasi eksplisit, dan strategi pembaruan struktur dependensi berbasis delta; (b) implementasi enam algoritma pembanding secara setara agar perbandingan pada tahap \textit{Evaluation} tidak bias oleh perbedaan kualitas implementasi; dan (c) konstruksi dataset ABAC-Lab dengan tingkat korelasi atribut dan \textit{drift} terkontrol, mengingat dataset semacam ini belum tersedia sebagai sumber daya siap pakai.
    \item \textit{Evaluation} -- Mengukur dan membandingkan hasil demonstrasi terhadap metrik (kualitas kebijakan/\textit{F-score} dan WSC\label{first:wsc}, jejak memori, serta waktu komputasi), disertai uji signifikansi statistik perbandingan antar-algoritma terhadap H1-H4.
    \item \textit{Communication} -- Mengomunikasikan permasalahan, artefak CoDA-EDA, dan hasil evaluasinya melalui buku tesis serta naskah publikasi ilmiah terkait.
\end{enumerate}

Meskipun keenam aktivitas disajikan berurutan, proses DSRM bersifat iteratif, sehingga penulis dimungkinkan kembali ke aktivitas \textit{Design and Development} apabila hasil \textit{Demonstration} dan \textit{Evaluation} menunjukkan kebutuhan penyempurnaan rancangan \parencite{peffers_design_2007}.